\section{Deep Learning for Ghost Imaging Reconstruction}
\label{sec:ghost_imaging}

The application of deep learning to ghost imaging reconstruction has been the most active area at the intersection of quantum imaging and neural networks. This section reviews the principal approaches, organised by network architecture.

\subsection{Fully Connected and CNN-Based Approaches}

\citet{lyu2017} provided the first demonstration that a neural network could reconstruct ghost images from bucket detector measurements. Their approach used a fully connected network trained on pairs of bucket signals and ground-truth images, achieving recognisable reconstruction of $28 \times 28$ handwritten digits at a sampling ratio as low as 12.5\%---a dramatic improvement over the correlation-based method (Eq.~\ref{eq:gi_correlation}), which requires sampling ratios approaching 100\% for comparable quality.

\citet{he2018} extended this work with a convolutional architecture, demonstrating improved generalisation to object categories not seen during training. Their CNN employed multiple convolutional layers with batch normalisation and ReLU activations, and was validated on both simulated and experimental ghost imaging data. Notably, the CNN-based approach maintained reasonable reconstruction quality even at sampling ratios below 10\%, where traditional methods produce severely degraded images.

\citet{shimobaba2018} applied deep learning specifically to computational ghost imaging, replacing the correlation computation of Eq.~\eqref{eq:gi_correlation} with a learned mapping. Their results confirmed that deep learning consistently outperforms traditional correlation-based reconstruction, particularly in low-measurement regimes.

\subsection{GAN-Based Reconstruction}

Generative adversarial networks \citep{goodfellow2014} have been applied to ghost imaging to improve the perceptual quality of reconstructions. The adversarial training framework encourages the generator to produce images that are statistically indistinguishable from real images, yielding reconstructions with sharper edges and more realistic textures than those obtained from MSE-trained networks.

However, GAN-based methods carry an inherent risk of \emph{hallucination}---generating plausible but incorrect image features---which is particularly concerning in scientific imaging applications where fidelity to the true object is paramount. This trade-off between perceptual quality and pixel-level accuracy remains an active area of investigation.

\subsection{End-to-End Learned Approaches}

A significant advance was made by \citet{wang2019}, who proposed jointly optimising the illumination patterns and the reconstruction network in an end-to-end differentiable framework. Rather than using random illumination patterns and learning only the reconstruction mapping, their approach learns \emph{optimal measurement strategies} tailored to the reconstruction network.

This end-to-end optimisation achieved 3--5\,dB improvement in PSNR over methods that separately optimise illumination and reconstruction, demonstrating that the measurement and reconstruction stages are fundamentally coupled and benefit from joint design. The approach also provides insight into what constitutes an ``informative'' measurement in the context of ghost imaging.

\subsection{Comparative Analysis}

Table~\ref{tab:gi_comparison} summarises the key characteristics and reported performance of the reviewed deep learning methods for ghost imaging reconstruction.

\begin{table*}[t]
\centering
\footnotesize
\caption{Comparison of deep learning methods for ghost imaging reconstruction.}
\label{tab:gi_comparison}
\begin{tabular}{@{}lp{2.2cm}cccl@{}}
\toprule
\textbf{Reference} & \textbf{Architecture} & \textbf{Size} & \textbf{Sampling} & \textbf{Data} & \textbf{Key Contribution} \\
\midrule
Lyu et al.\ [2017]       & Fully connected  & $28\times28$   & 12.5\%  & Sim.     & First DL-based ghost imaging reconstruction \\
He et al.\ [2018]        & CNN              & $28\times28$   & $<$10\% & Sim.+Exp. & Improved generalisation to unseen categories \\
Shimobaba et al.\ [2018] & CNN              & $64\times64$   & 25\%    & Sim.     & Applied DL to computational GI \\
Wang et al.\ [2019]      & End-to-end CNN   & $64\times64$   & 6.25\%  & Sim.     & Joint optimisation of patterns and network \\
\bottomrule
\end{tabular}
\end{table*}

Several trends emerge from this comparison. First, all methods achieve usable reconstruction quality at sampling ratios well below the Nyquist limit, confirming the strong regularising effect of learned image priors. Second, the field has progressed from simple fully connected architectures to more sophisticated end-to-end frameworks. Third, there is a notable lack of standardised benchmarks: different papers use different image sizes, datasets, noise models, and evaluation metrics, making rigorous cross-method comparison difficult.

\subsection{Discussion}

The reviewed works establish that deep learning is a transformative tool for ghost imaging reconstruction. However, several limitations persist. Most studies operate on small image sizes ($28\times28$ to $64\times64$); scalability to high-resolution imaging remains underexplored. The reliance on simulated data in many studies raises questions about robustness to real experimental noise and detector imperfections. Finally, the fundamental question of whether quantum correlations provide an advantage over classical correlations \emph{when combined with deep learning} has not been systematically addressed.
